% -----------------------------------------------------------------------
%  lifelines-hc Application Note — Bioinformatics journal
%
%  Class:   oup-authoring-template  (OUP standard; pre-installed on
%           Overleaf and available via CTAN / TeX Live)
%  Style:   contemporary, large, unnumbered section heads, numbered refs
%
%  Compile: pdflatex main && bibtex main && pdflatex main && pdflatex main
%  Overleaf: upload all files; template class is already available there.
% -----------------------------------------------------------------------
\documentclass[unnumsec,webpdf,contemporary,large]{oup-authoring-template}

% ---- packages ----
\usepackage{microtype}
\usepackage{booktabs}
\usepackage{listings}
\usepackage{xcolor}

\lstset{%
  language=Python,
  basicstyle=\ttfamily\small,
  keywordstyle=\color{blue!65!black}\bfseries,
  commentstyle=\color{gray!70},
  stringstyle=\color{green!45!black},
  breaklines=true,
  frame=tb,
  framesep=4pt,
  backgroundcolor=\color{gray!7},
  xleftmargin=1.2em,
  xrightmargin=1.2em,
  aboveskip=0.8em,
  belowskip=0.5em,
  showstringspaces=false,
}

\graphicspath{{../figs/}}

% -----------------------------------------------------------------------
\begin{document}

% ---- journal metadata (leave most blank; editor fills on acceptance) ----
\journaltitle{Biometrika}
\DOI{10.1093/biomet/asaf075}
\copyrightyear{2025}
\pubyear{2026}
\access{Advance Access Publication Date: Day Month Year}
\appnotes{Application Note}

\firstpage{1}

% ---- title ----
\title[{\texttt{lifelines-hc}}: Higher Criticism survival testing]%
      {{\texttt{lifelines-hc}}: a Python package for Higher Criticism
       testing in two-sample survival analysis under non-proportional
       hazards}

% ---- authors ----
\author[1,$\ast$]{Alon Kipnis}
\author[2]{Ben Galili}
\author[1,2]{Zohar Yakhini}

\authormark{Kipnis}

\address[1]{%
  \orgdiv{School of Computer Science},
  \orgname{Reichman University},
  \orgaddress{%
    \street{Herzliya},
    \postcode{4610101},
    \country{Israel}%
  }%
}

\address[2]{%
  \orgdiv{Department of Computer Science},
  \orgname{Technion -- Israel Institute of Technology},
  \orgaddress{%
    \street{Haifa},
    \postcode{3200003},
    \country{Israel}%
  }%
}

\corresp[$\ast$]{%
  E-mail: \href{mailto:alon.kipnis@runi.ac.il}{alon.kipnis@runi.ac.il}%
}

% Dates — leave blank; editor fills on acceptance (month arg must be numeric)
\received{}{0}{}
\revised{}{0}{}
\accepted{}{0}{}

%\editor{Associate Editor: Name}

% ---- structured abstract ----
% Bioinformatics Application Notes: Motivation | Results | Availability
% Use \\ between subsections; no blank lines inside \abstract{...}.
\abstract{%
\textbf{Motivation:}
The log-rank test is a standard tool for two-sample survival comparison and has good power against the widely studied proportional-hazards alternatives. However, \citet{kipnis_biometrika} showed that it can lose power against sparse hazard departures, where the hazard difference is concentrated in a small number of time intervals whose locations along the follow-up are unknown a priori. Weighted log-rank tests, such as Gehan-Wilcoxon, Tarone-Ware, Peto-Prentice, and Fleming-Harrington, each impose a pre-specified temporal emphasis, and are therefore also insensitive to sparse departures occurring outside their emphasized time regions.
%
\textbf{Results:}
We introduce \texttt{lifelines-hc}, a Python package extending the
\texttt{lifelines} survival library with the HCHG test: Higher Criticism
(HC) applied to per-interval HyperGeometric (HG) p-values from
\citet{kipnis_biometrika}. 
HCHG is a test for right-censored two-sample survival data. It detects sparse hazard departures at unknown locations without committing to any temporal pattern. Across four biological case studies
— an immuno-oncology trial (CheckMate~057 PFS, $n=582$), a targeted
kinase-inhibitor trial in metastatic prostate cancer (COMET-1 OS,
$n=1028$), an adjuvant bisphosphonate trial (AZURE DFS, $n=3359$),
and the 
Copenhagen Study group for Liver diseases (CSL) liver-cirrhosis trial with fully public individual patient data ($n=446$) — HCHG achieves $p \leq 0.014$ in every dataset, while the log-rank test is non-significant throughout ($p \geq 0.26$); among the weighted log-rank tests, only Gehan-Wilcoxon and Peto-Prentice on CheckMate~057 reach $p < 0.05$.\\
%
\textbf{Availability and Implementation:}
\texttt{lifelines-hc} is freely available under the MIT licence at
\url{https://github.com/alonkipnis/lifelines-hc} and via
\texttt{pip install lifelines-hc}; it requires Python~${\geq}\,3.8$ and
\texttt{lifelines}~${\geq}\,0.29$. Full analysis scripts and per-method
results tables are provided as Supplementary Data.%
}

\keywords{survival analysis; higher criticism; sparse signal detection;
  hypergeometric test; hypothesis testing; Python}

\maketitle


% =======================================================================
\section{Introduction}
% =======================================================================

The Kaplan-Meier estimator \citep{kaplan1958} and the log-rank test
\citep{mantel1966} are the workhorses of two-sample survival analysis in
biomedical research. The log-rank statistic has good power against a wide
range of useful alternatives, particularly proportional-hazards (PH)
alternatives \citep{schoenfeld1981}. Nevertheless,
\citet{kipnis_biometrika} showed that the log-rank test is not optimal
against \emph{sparse} hazard departures, in which the hazard difference is
concentrated in a small, unknown subset of time intervals.

Such localized or non-proportional departures arise naturally in several
biomedical settings. Immune checkpoint inhibitors may require an
immune-priming phase before producing durable tumour suppression, leading
to delayed treatment effects and crossing survival curves; for example,
in the second-line nonsquamous NSCLC trial CheckMate~057, nivolumab
improved overall survival relative to docetaxel, while the
progression-free-survival curves crossed, consistent with delayed benefit
\citep{borghaei2015}. Targeted therapies can also produce endpoint-specific
benefits over a limited follow-up window that are diluted in global
survival summaries: in the COMET-1 mCRPC trial, cabozantinib markedly
improved bone-scan response at week~12 and radiographic
progression-free survival, but did not significantly improve overall
survival \citep{smith2016}. Subgroup-dependent biology can similarly
generate composite, time-varying hazard differences: in the AZURE trial
of adjuvant zoledronic acid, no overall benefit was observed in the full
population, but benefit was seen among women with established menopause,
consistent with an effect mediated by the postmenopausal bone
microenvironment \citep{coleman2011,coleman2014}.
Time-varying treatment effects can likewise produce multi-window hazard
departures: in the CSL liver-cirrhosis trial, prednisone was harmful early
(infection and fluid retention) but beneficial later through
anti-inflammatory action, yielding survival curves that cross more than once
\citep{csl_data}. This canonical non-proportional-hazards dataset is widely
used to illustrate time-varying-coefficient models and provides fully public
individual patient data via the CRAN \texttt{timereg} package
\citep{csl_data} and the SurvSet repository \citep{survset2022}.

The need for log-rank alternatives in survival analysis is well recognised,
and several weighted log-rank tests have been developed for non-proportional-hazards
settings \citep{gehan1965,peto1972,taroneware1977,flemingharrington1991}.
Each applies a temporal weight function, emphasising early, intermediate,
or late events. Each of these tests achieves good power when the signal matches the intended temporal pattern. However, choosing an appropriate weight requires knowing in advance \emph{where} the hazard departure will occur. When the departure is sparse but its location is unknown a priori, all weighted log-rank tests can miss it simultaneously, as we demonstrate below.

\citet{kipnis_biometrika} proposed the HCHG test — the Higher Criticism
(HC) statistic of \citet{donoho2004} applied to per-interval
HyperGeometric (HG) p-values — specifically for this setting. HCHG
divides the follow-up into equal-width intervals, derives a hypergeometric
p-value for the event-count comparison in each interval under the null of
equal group-specific hazards, and applies the HC statistic to test whether
the body of these p-values is collectively implausible
under the global null. Because HCHG does not pre-specify any temporal
weight, it retains power for sparse departures wherever they occur along
the follow-up axis.

We present \texttt{lifelines-hc}, a Python package that implements HCHG and integrates with the popular \texttt{lifelines} library \citep{davidsonpilon2019}.


% =======================================================================
\section{The \texttt{lifelines-hc} package}
% =======================================================================

\subsection*{The HCHG statistic}

HCHG \citep{kipnis_biometrika} combines two components: hypergeometric
interval p-values and the HC statistic of \citet{donoho2004,donoho2015special}.

\textbf{Hypergeometric interval p-values.}
Given two groups with right-censored observations $(T_i, E_i, G_i)$,
partition the follow-up range into $K$ equal-width intervals. For
interval $k$, let $n_k$ be the total events, $r_k^A$ and $r_k^B$ the
subjects at risk in groups A and B, and $x_k$ the events in group B.
Under the null of equal group-specific hazards,
%
\begin{equation}\label{eq:hypergeom}
  x_k \;\Big|\; n_k,\, r_k^A,\, r_k^B
  \;\sim\;
  \mathrm{Hypergeometric}\!\left(r_k^A + r_k^B,\; r_k^B,\; n_k\right).
\end{equation}
%
yielding a one-sided p-value $p_k$ for each interval using an exact test.

\textbf{HC aggregation.}
Ordering the $K$ p-values as $p_{(1)} \leq \cdots \leq p_{(K)}$, the
HC statistic \citep{donoho2004,donoho2015special} is applied to this collection:
%
\begin{equation}\label{eq:hc}
  \mathrm{HCHG}_K
  \;=\;
  \max_{1 \leq j \leq \lfloor\gamma K\rfloor}
  \frac{j/K - p_{(j)}}{\sqrt{p_{(j)}\!\left(1 - p_{(j)}\right)/K}},
\end{equation}
%
where $\gamma \in (0,1)$ (default $0.2$) restricts scanning to the most
extreme fraction of interval p-values \citep{donoho2004,donoho2015special}. A global
p-value is obtained by permuting group labels
($n_\mathrm{perms} = 2000$ in the analyses reported here). For the two-sided alternative,
the minimum of the two one-sided interval p-values replaces $p_k$
in~(\ref{eq:hc}).

\subsection*{Software design and API}

\texttt{lifelines-hc} exposes a function-based API consistent with
\texttt{lifelines.statistics}, accepting NumPy arrays of event times and
indicators:

\begin{lstlisting}
from lifelines_hc import (higher_criticism_test,
                           fisher_combination_test,
                           min_p_test,
                           suspected_deviations)

result = higher_criticism_test(
    T_A, T_B,
    event_observed_A=E_A, event_observed_B=E_B,
    n_intervals_to_pool=100, n_permutations=2000,
    alternative='both', gamma=0.2, seed=42)

print(result.p_value)       # permutation-calibrated p-value
print(result.test_statistic)  # HCHG statistic value
\end{lstlisting}

\noindent
The companion \texttt{suspected\_deviations()} function returns a
\texttt{pandas.DataFrame} of per-interval p-values with a
\texttt{suspected} boolean column marking HCHG-flagged intervals,
enabling overlay of diagnostic shading on Kaplan-Meier plots. Two
complementary omnibus tests are also provided: Fisher combination
\citep{fisher1932}, which aggregates evidence via $-2\sum_k \log p_k$,
and a Bonferroni-corrected MinP test. All three share the same
permutation infrastructure and the same hypergeometric interval p-values.
The package additionally ships with plotting utilities
(\texttt{plot\_km\_with\_hc}, \texttt{plot\_pvalue\_profile}) and a
convenience wrapper \texttt{run\_all\_tests()} that benchmarks HCHG,
Fisher, MinP, log-rank, and the four weighted alternatives in a single
call, returning a \texttt{pandas.DataFrame} for easy comparison.


% =======================================================================
\section{Results}
% =======================================================================

We demonstrate \texttt{lifelines-hc} on four biological datasets with
distinct non proportional-hazards (PH) structures (Figure~\ref{fig:main}), benchmarking HCHG
against the log-rank test and four weighted alternatives.
Individual patient data for the first three trials were reconstructed from
published Kaplan-Meier curves; the CSL trial provides publicly available individual patient data.
Full per-method p-value tables are provided in Supplementary Data.

\subsection*{Immuno-oncology: crossing survival curves}

CheckMate~057 \citep{borghaei2015} randomised 582 patients with advanced
non-squamous NSCLC to nivolumab ($n=292$) or docetaxel ($n=290$).
Progression-free survival showed crossing curves: docetaxel provided a
short-term advantage during immune priming, subsequently overtaken by the
durable nivolumab response (Figure~\ref{fig:main}A). Individual patient
data were reconstructed from the published Kaplan-Meier curve via the
algorithm of \citet{guyot2012} as implemented in the \texttt{kmdata}
R package.

The log-rank test is non-significant ($p = 0.351$) as the early excess
and late benefit partially cancel. Among the weighted alternatives,
Gehan-Wilcoxon ($p = 0.041$) and Peto-Prentice ($p = 0.049$) reach
borderline significance by detecting the early crossing, but
Tarone-Ware ($p = 0.358$) and Fleming-Harrington\,(1,1) ($p = 0.256$)
do not. HCHG achieves $p = 0.001$ and characterises the full crossing
structure with far greater power (Figure~\ref{fig:main}A,E) by flagging
both the early divergence and late separation windows.

\subsection*{Targeted therapy in prostate cancer: a transient early benefit}

The COMET-1 trial \citep{smith2016} randomised 1028 patients with
chemotherapy-pretreated mCRPC to cabozantinib ($n=682$) or prednisone
($n=346$) in a 2:1 ratio. Overall survival showed no significant
log-rank difference ($p = 0.262$). Cabozantinib, by inhibiting
MET and VEGFR2 kinases, significantly improved bone scan response at
week~12 (a pre-specified secondary endpoint), confirming genuine
biological activity against bone-metastatic disease.
This short-term activity did not translate into sustained overall survival (OS) benefit,
however, as resistance through alternative survival pathways rapidly
emerged. The result is a temporally concentrated early OS advantage
(months~3--7) that is subsequently lost (Figure~\ref{fig:main}B,F).

All four weighted alternatives are non-significant ($p \geq 0.19$):
the early-weighted Gehan-Wilcoxon ($p = 0.193$) and Peto-Prentice
($p = 0.206$) come closest but are diluted by the null period after
the bone-response window. HCHG ($p = 0.014$) aggregates evidence from the specific early intervals
where the transient effect is concentrated.

\subsection*{Adjuvant bisphosphonate therapy: a menopause-dependent effect}

The AZURE trial \citep{coleman2011,coleman2014} randomised 3359 early-stage
breast cancer patients to standard adjuvant therapy with ($n=1681$) or
without ($n=1678$) zoledronic acid. Disease-free survival showed no
significant log-rank difference ($p = 0.305$) because the drug's
bone-microenvironment benefit depends on the low-oestrogen state of
postmenopause, which was absent in the many premenopausal patients
enrolled at baseline. As these patients progressively transitioned to
postmenopause during the decade-long follow-up, the hazard difference
accumulated in a specific mid-to-late window (months 20--60) rather than
uniformly (Figure~\ref{fig:main}C,G).

All four weighted alternatives are non-significant ($p \geq 0.28$):
no single temporal weight captures a pattern distributed across this
window while remaining unaffected by the null early period. HCHG
($p = 0.005$) detects this scattered but collectively significant signal.
The result is validated by published subgroup analyses showing significant disease-free survival (DFS) benefit in postmenopausal patients \citep{coleman2011,coleman2014},
confirming that HCHG is detecting a genuine biological effect.

\subsection*{Corticosteroid therapy in liver cirrhosis: a time-varying effect}

The CSL trial \citep{csl_data} randomised 446 patients with histologically
verified liver cirrhosis to prednisone ($n=226$) or placebo ($n=220$) in
Copenhagen (1962--69), with 270 deaths over follow-up of up to
$\sim$13 years. Unlike the preceding case studies, the raw individual
patient data are genuinely public---distributed in the CRAN \texttt{timereg}
package (dataset \texttt{csl}) and the open SurvSet repository
\citep{survset2022}---rather than reconstructed from a published
Kaplan-Meier curve.

Corticosteroids exert a time-varying effect in cirrhosis: early harm through
immunosuppression, infection, and fluid retention is followed by later
anti-inflammatory benefit in actively inflamed subgroups. The survival
curves cross more than once, with the treatment effect concentrated in a
few non-contiguous windows (Figure~\ref{fig:main}D,H).

The log-rank test is non-significant ($p = 0.384$), and all four weighted
alternatives are as well ($p \geq 0.13$): Gehan-Wilcoxon ($p = 0.51$),
Tarone-Ware ($p = 0.38$), Peto-Prentice ($p = 0.45$), and
Fleming-Harrington\,(1,1) ($p = 0.14$). No single temporal weight
captures a multi-window, sign-changing departure. HCHG achieves
$p = 0.002$ and Fisher combination $p = 0.034$, demonstrating detection
of sparse non-proportional-hazards departures on genuinely public raw data.

Across the four case studies, HCHG achieves $p \leq 0.014$ in every domain,
while the log-rank test is non-significant throughout ($p \geq 0.26$). The
two borderline significances for Gehan-Wilcoxon ($p = 0.041$) and
Peto-Prentice ($p = 0.049$) in CheckMate~057 are the only instances where
any weighted alternative reaches $p < 0.05$; elsewhere the weighted tests
return $p \geq 0.13$. These reflect their early emphasis coincidentally
aligning with the crossing-curve signal rather than general sensitivity
to non-PH. The four patterns differ---early
and late crossing, transient early benefit, mid-to-late accumulated effect,
and multi-window sign-changing departure---yet HCHG detects all without
per-dataset tuning. Fisher combination is also significant in all four
domains, consistent with sparse localised signals.


% =======================================================================
\section{Conclusion}
% =======================================================================

\texttt{lifelines-hc} provides an accessible Python implementation of
HCHG for two-sample right-censored survival data. HCHG is specifically
designed for sparse hazard departures at unknown temporal locations:
unlike weighted log-rank alternatives, it requires no a priori
specification of where the hazard difference will occur and retains
power for any concentrated departure along the follow-up axis. The
package integrates transparently with \texttt{lifelines} workflows,
returns diagnostic interval-level information identifying which time
windows drive the signal, and includes Fisher combination and MinP tests
as complementary omnibus alternatives sharing the same hypergeometric
infrastructure.


% =======================================================================
%  Back matter
% =======================================================================

% \section*{Acknowledgements}
% TODO.

\section*{Conflict of Interest}
None declared.

\section*{Funding}
The work of Alon Kipnis was supported in parts by the US-Israel Binational Science Foundation (BSF) grant 2022124.

\section*{Data availability}
Individual patient data for CheckMate~057, COMET-1, and AZURE were
reconstructed from published Kaplan-Meier curves using the algorithm of
\citet{guyot2012} as implemented in the \texttt{kmdata} R package
(\url{https://github.com/raredd/kmdata}). CSL liver-cirrhosis individual
patient data are fully public via the CRAN \texttt{timereg} package
(dataset \texttt{csl}) and the SurvSet repository \citep{survset2022}
(\url{https://github.com/ErikinBC/SurvSet}).
All analysis code is available at
\url{https://github.com/alonkipnis/lifelines-hc/tree/main/AppNote}.


% =======================================================================
%  Figure — full-width across both columns
% =======================================================================

\begin{figure*}[t!]
  \centering
  \includegraphics[width=\textwidth]{figure1_4dom}
  \caption{%
    \textbf{Higher Criticism detects non-proportional hazard departures
    across four clinical datasets with distinct temporal patterns.}
    Top row (A--D): Kaplan-Meier curves with HCHG-flagged intervals
    (shaded). Bottom row (E--H): per-interval \emph{signed}
    $-\!\log_{10}$(hypergeometric p-value). Upward bars mark excess
    events in the treatment arm, downward bars excess in the control arm.
    Dashed lines are the per-direction HCHG thresholds; bars reaching
    beyond either line are highlighted in red and coincide exactly with
    the intervals shaded above.
    %
    \textbf{(A,E)}~CheckMate~057 PFS (nivolumab vs.\ docetaxel, $n=582$);
    log-rank $p=0.351$, HCHG $p=0.001$. Curves cross at $\sim\!3$ months
    before nivolumab gains a durable late advantage.
    %
    \textbf{(B,F)}~COMET-1 OS (cabozantinib vs.\ prednisone, $n=1028$);
    log-rank $p=0.262$, HCHG $p=0.014$. Early OS advantage (months~3--7)
    during bone-lesion response, fading as resistance develops.
    %
    \textbf{(C,G)}~AZURE DFS (zoledronic acid vs.\ control, $n=3359$);
    log-rank $p=0.305$, HCHG $p=0.005$. Menopause-dependent benefit
    accumulates in months~20--60.
    %
    \textbf{(D,H)}~CSL OS (prednisone vs.\ placebo, liver cirrhosis,
    $n=446$); log-rank $p=0.384$, HCHG $p=0.002$. Time-varying
    corticosteroid effect with multiple curve crossings; \emph{fully
    public} raw individual patient data (CRAN \texttt{timereg} / SurvSet).
    In every dataset the log-rank test is non-significant while HCHG is
    significant. For the first three, individual patient data
    reconstructed via \citet{guyot2012}.%
  }
  \label{fig:main}
\end{figure*}


% =======================================================================
%  Bibliography
% =======================================================================

% oup-authoring-template.bst is available on Overleaf; locally use plainnat.
% Switch to \bibliographystyle{oup-authoring-template} when uploading to Overleaf.
\bibliographystyle{plainnat}
\bibliography{refs}

\end{document}
